% !TEX root = ../../ctfp-print.tex

\lettrine[lhang=0.17]{プ}{ログラミングから} 離れて、ガチガチの数学に飛び込んでいるように感じるかもしれません。しかし、プログラミングにおける次の大革命が何をもたらすか、またそれを理解するためにどんな数学が必要になるかは誰にも分かりません。連続時間を扱う関数型リアクティブプログラミング、依存型によるHaskellの型システムの拡張、あるいはプログラミングにおけるホモトピー型理論の探求など、非常に興味深いアイデアが世の中に出回っています。

これまで私は、型を値の\emph{集合}と気軽に同一視してきました。これは厳密には正しくありません。なぜなら、そのようなアプローチは、プログラミングにおいて値を\emph{計算する}という事実を考慮していないからです。計算とは時間のかかるプロセスであり、極端な場合には終了しないこともあります。発散する計算は、チューリング完全なすべての言語に存在します。

集合論が計算機科学あるいは数学そのものの基礎として最適ではない理由は、基礎的なレベルでも存在します。良い例えは、集合論が特定のアーキテクチャに縛られたアセンブリ言語であるというものです。異なるアーキテクチャで数学を動かしたければ、より一般的なツールを使う必要があります。

一つの可能性は、集合の代わりに空間を使うことです。空間はより多くの構造を持ち、集合を使わずに定義できます。空間に通常関連付けられるのが位相幾何学であり、これは連続性などの概念を定義するために必要です。そして位相幾何学への従来のアプローチは、お察しの通り、集合論を通したものです。特に、部分集合の概念は位相幾何学の中心にあります。当然ながら、圏論者はこのアイデアを$\Set$以外の圏に一般化しました。集合論の代替として機能するのに適切な性質を持つ圏の型を\newterm{トポス}（複数形：トポイ）と呼び、それは特に部分集合の一般化された概念を提供します。

\section{部分対象分類子}

要素を使う代わりに関数を使って部分集合の概念を表現することから始めましょう。ある集合$a$から$b$への任意の関数$f$は、$b$の部分集合——$a$の$f$による像——を定義します。しかし、同じ部分集合を定義する関数は多数あります。もっと具体的にする必要があります。まず、単射的な関数——複数の要素を一つに押し込めない関数——に注目するかもしれません。単射関数は一つの集合を別の集合に「注入」します。有限集合の場合、単射関数は一方の集合の要素から他方の集合の要素へと向かう並行した矢印として視覚化できます。もちろん、最初の集合が2番目の集合より大きいと矢印が必ず収束してしまいます。まだいくらかの曖昧さが残っています。別の集合$a'$と、その集合から$b$への別の単射関数$f'$が、同じ部分集合を選ぶことがあるかもしれません。しかし、そのような集合は$a$と同型でなければならないことは容易に確認できます。この事実を使って、部分集合をドメインの同型によって関係付けられた単射関数の族として定義できます。より正確には、二つの単射関数：
\begin{align*}
  f  & \Colon a \to b  \\
  f' & \Colon a' \to b
\end{align*}
が同値であるとは、同型射：
\[h \Colon a \to a'\]
が存在して：
\[f = f' \circ h\]
を満たすことです。このような同値な単射の族が$b$の部分集合を定義します。

\begin{figure}[H]
  \centering
  \includegraphics[width=0.4\textwidth]{images/subsetinjection.jpg}
\end{figure}

\noindent
この定義は、単射関数をモノ射に置き換えることで、任意の圏に持ち上げられます。念のため言えば、$a$から$b$へのモノ射$m$は普遍性によって定義されます。任意の対象$c$と任意の射の対：
\begin{align*}
  g  & \Colon c \to a \\
  g' & \Colon c \to a
\end{align*}
に対して：
\[m \circ g = m \circ g'\]
であれば$g = g'$でなければなりません。

\begin{figure}[H]
  \centering
  \includegraphics[width=0.4\textwidth]{images/monomorphism.jpg}
\end{figure}

\noindent
集合においては、この定義は関数$m$がモノ射で\emph{ない}とはどういうことかを考えると理解しやすいです。それは$a$の2つの異なる要素を$b$の一つの要素に写すことになります。その場合、それら2つの要素でのみ異なる関数$g$と$g'$を見つけられます。$m$との後合成でその違いが隠されてしまいます。

\begin{figure}[H]
  \centering
  \includegraphics[width=0.4\textwidth]{images/notmono.jpg}
\end{figure}

\noindent
部分集合を定義する別の方法があります。特性関数と呼ばれる一つの関数を使う方法です。これは集合$b$から2要素集合$\Omega$への関数$\chi$です。この集合の一方の要素が「真」、もう一方が「偽」として指定されます。この関数は、$b$の要素のうち部分集合のメンバーであるものに「真」を、そうでないものに「偽」を割り当てます。

$\Omega$の要素を「真」として指定するとはどういう意味かを特定する必要があります。標準的なトリックを使いましょう。一点集合から$\Omega$への関数を使います。この関数を$\mathit{true}$と呼びます：
\[\mathit{true} \Colon 1 \to \Omega\]

\begin{figure}[H]
  \centering
  \includegraphics[width=0.4\textwidth]{images/true.jpg}
\end{figure}

\noindent
これらの定義は、部分対象とは何かを定義するだけでなく、要素について語ることなく特別な対象$\Omega$を定義するような形で組み合わせることができます。アイデアは、射$\mathit{true}$が「一般的な」部分対象を表すようにしたいということです。$\Set$においては、これは2要素集合$\Omega$から一要素の部分集合を選びます。これはできる限り一般的なものです。$\Omega$にはその部分集合に\emph{含まれない}要素がもう一つあるので、これは明らかに真の部分集合です。

より一般的な設定では、$\mathit{true}$を終対象から\emph{分類対象}$\Omega$へのモノ射として定義します。\emph{部分対象分類子}は$\Omega$と$\mathit{true}$を合わせたものからなります。しかし、分類対象を定義する必要があります。この対象と特性関数を結びつける普遍性が必要です。$\Set$においては、特性関数$\chi$に沿った$\mathit{true}$の引き戻しが、部分集合$a$とそれを$b$に埋め込む単射関数の両方を定義することが分かります。引き戻し図式はこのようになります：

\begin{figure}[H]
  \centering
  \includegraphics[width=0.4\textwidth]{images/pullback.jpg}
\end{figure}

\noindent
この図式を分析しましょう。引き戻しの方程式は：
\[\mathit{true} \circ \mathit{unit} = \chi \circ f\]
関数$\mathit{true} \circ \mathit{unit}$は$a$のすべての要素を「真」に写します。したがって$f$は$a$のすべての要素を、$\chi$が「真」である$b$の要素に写さなければなりません。それらは定義上、特性関数$\chi$によって指定される部分集合の要素です。$f$の像はまさに問題の部分集合です。引き戻しの普遍性は$f$が単射であることを保証します。

この引き戻し図式は$\Set$以外の圏で分類対象を定義するために使えます。そのような圏は終対象を持つ必要があり、それによって$\mathit{true}$というモノ射を定義できます。また引き戻しも持つ必要があります——実際の要件はすべての有限極限を持つことです（引き戻しは有限極限の一例です）。これらの仮定のもと、すべてのモノ射$f$に対して、引き戻し図式を完成させる唯一の射$\chi$が存在するという性質によって分類対象$\Omega$を定義します。

最後の文を分析しましょう。引き戻しを構成するとき、3つの対象$\Omega$、$b$、$1$と、2つの射$\mathit{true}$と$\chi$が与えられます。引き戻しの存在は、図式を可換にする2つの射$f$と$\mathit{unit}$（後者は終対象の定義によって一意に決まります）を備えた最良の対象$a$を見つけられることを意味します。

ここでは異なる方程式系を解いています。$a$と$b$の両方を変化させながら、$\Omega$と$\mathit{true}$を解いています。与えられた$a$と$b$に対して、モノ射$f \Colon a \to b$が存在する場合もあれば、存在しない場合もあります。しかし存在する場合は、それがある$\chi$の引き戻しであってほしいのです。さらに、この$\chi$が$f$によって一意に決まることを望んでいます。

モノ射$f$と特性関数$\chi$の間に一対一対応があるとは言えません。なぜなら引き戻しは同型を除いてのみ一意だからです。しかし、同値な単射の族としての部分集合の先ほどの定義を思い出してください。$b$の部分対象を$b$への同値なモノ射の族として定義することで一般化できます。このモノ射の族は、図式の同値な引き戻しの族と一対一対応にあります。

したがって、$b$の部分対象の集合$\mathit{Sub}(b)$をモノ射の族として定義し、それが$b$から$\Omega$への射の集合と同型であることが分かります：
\[\mathit{Sub}(b) \cong \cat{C}(b, \Omega)\]
これは2つの関手の自然同型になっています。言い換えれば、$\mathit{Sub}(-)$は表現可能（反変）関手であり、その表現対象が$\Omega$です。

\section{トポス}

トポスとは次の性質を持つ圏です：

\begin{enumerate}
  \tightlist
  \item
        デカルト閉である：すべての積、終対象、および指数対象（積への右随伴として定義）を持つ。
  \item
        すべての有限図式に対して極限を持つ。
  \item
        部分対象分類子$\Omega$を持つ。
\end{enumerate}

この性質の集合により、トポスはほとんどの応用において$\Set$の代替として申し分ないものになります。また、定義から導かれる追加の性質もあります。例えば、トポスは始対象を含むすべての有限余極限を持ちます。

部分対象分類子を終対象の2つのコピーの余積（和）として定義したいと思うかもしれません——$\Set$ではそうです——しかし、より一般化したいと考えています。これが成り立つトポスはブール型と呼ばれます。

\section{トポスと論理}

集合論では、特性関数は集合の要素の性質を定義するものとして解釈できます——一部の要素に対しては真、他の要素に対しては偽である\newterm{述語}です。述語$\mathit{isEven}$は自然数の集合から偶数の部分集合を選び出します。トポスでは、述語の概念を対象$a$から$\Omega$への射として一般化できます。これが$\Omega$が真理対象と呼ばれることがある理由です。

述語は論理の構成要素です。トポスは論理を研究するために必要なすべての道具立てを含んでいます。論理的連言（論理\emph{積}）に対応する積、選言（論理\emph{和}）のための余積、そして含意のための指数対象を持ちます。排中律（あるいは同値な二重否定の除去）を除いて、論理のすべての標準的な公理がトポスで成立します。これがトポスの論理が構成主義的あるいは直観主義的論理に対応する理由です。

直観主義論理は着実に支持を広げており、計算機科学から予期せぬ支持を得ています。排中律の古典的な概念は絶対的な真理の存在という信念に基づいています：どんな命題も真か偽かのどちらかであり、あるいは古代ローマ人が言うように、\emph{tertium non datur}（第三の選択肢はない）です。しかし、あることが真か偽かを知る唯一の方法は、それを証明または反証できる場合です。証明とはプロセスであり、計算です——そして計算には時間とリソースがかかることが分かっています。場合によっては、決して終了しないかもしれません。有限の時間内に証明できない命題が真であると主張するのは意味をなしません。より微妙な真理対象を持つトポスは、興味深い論理をモデル化するための、より一般的なフレームワークを提供します。

\section{課題}

\begin{enumerate}
  \tightlist
  \item
        特性関数に沿った$\mathit{true}$の引き戻しである関数$f$が単射でなければならないことを示せ。
\end{enumerate}
